\scanpage{251}\mspage{- \handwritten{2} -}
n$^{\mathrm{o}}$ \handwritten{293}

un corps des fractions de $\mathrm o/\mathrm p$ ; soit $R(\theta_1,..,\theta_n)$ une extension
transcendante pure de $R$ engendrée par $n$ éléments $\theta_1,..,\theta_n$ algébriquement
indépendants par rapport à $R$. L'homomorphisme canonique de
$\mathrm o$ sur $\mathrm o/\mathrm p$ se prolonge en un homomorphisme $\omega$ de $\mathrm o[t_1,..,t_n]$
dans $R[\theta_1,..,\theta_n]$ qui applique $t_i$ sur $\theta_i$ $(1\leq i\leq n)$ ; le noyau
de cet homomorphisme est $\mathrm{po}[t_1,..,t_n]$, comme il résulte tout
de suite du fait que $\theta_1,..,\theta_n$ sont algébriquement indépendants
sur $R$ ; par ailleurs, $\omega(\mathrm o[t_1,..,t_n])$ est un domaine d'intégrité
ce qui montre bien que $\mathrm{po}[t_1,..,t_n]$ est premier. De plus
$\omega(\mathrm o[t_1,..,T_n])$ n'est pas un corps. En effet, dans le cas contraire,
ce corps devrait contenir $R$, donc être identique à $R[\theta_1,..,\theta_n]$
or ce dernier anneau n'est pas un corps puisque $n>0$. O\deleted{\unsure{u|r}}n en conclut
que $\mathrm{po}[t_1,..,t_n]$ n'est pas un idéal premier maximal de
$\mathrm o[t_1,..,t_n]$. Par contre, $\mathrm q$ étant un idéal premier maximal de
$\mathrm r$, et $\mathrm r$ étant entier sur $\mathrm o[t_1,\ldots,t_n]$, $\mathrm q\cap\mathrm o[t_1,\ldots,t_n]$
est un idéal premier maximal de $\mathrm o[t_1,..,t_n]$ ; cet idéal contient
$\mathrm{po}[t_1,\ldots,t_n]$, et en est distinct en vertu de ce qu'on vient
de dire. Comme $\mathrm o[t_1,\ldots,t_n]$ est intégralement clos, on sait
qu'il y a un idéal premier $\mathrm n$ de $\mathrm r$ tel que $\mathrm n\subset\mathrm q$, $\mathrm n\cap\mathrm o[t_1,..,t_n]$
$=\mathrm{po}[t_1,\ldots,t_n]$ ; cet idéal est contenu dans $\mathrm q$, distinct de
$\mathrm q$ et contient $\mathrm p$.

\uline{Lemme} \handwritten{2} - \uline{Soient} $\mathrm o$ \uline{un domaine d'intégrité local}, $\mathrm m$ \uline{son idéal
premier maximal,} $t$ \uline{un élément d'un sur-corps du corps des fractions
de} $\mathrm o$ ; \uline{on suppose que} $\mathrm o$ \uline{est intégralement fermé dans} $\mathrm o[t]$,
\uline{que} $t\notin\mathrm o$, \uline{et que l'idéal engendré par} $\mathrm m$ \uline{dans} $\mathrm o[t]$ \uline{ne
contient pas} $1$. \uline{Si} $F$ \uline{est un polynôme à coefficients dans} $\mathrm o$ \uline{tel}
% archive sheet number: 250
